\documentclass[12pt,a4paper]{article}

\usepackage[utf8]{inputenc}
\usepackage[T1]{fontenc}
\usepackage[brazil]{babel}
\usepackage{lmodern}
\usepackage{geometry}
\usepackage{booktabs}
\usepackage{amsmath}
\usepackage{amssymb}
\usepackage{hyperref}
\usepackage{xcolor}
\usepackage{enumitem}

\geometry{margin=2.5cm}
\setlength{\parindent}{1.25cm}
\setlength{\parskip}{0.25em}
\hypersetup{
    colorlinks=true,
    linkcolor=blue!45!black,
    urlcolor=blue!45!black,
    citecolor=blue!45!black
}

\title{Agrupamento de Dados com BSAS, K-Means, Janelas de Parzen e KNN}
\author{Caique}
\date{\today}

\begin{document}

\maketitle

\begin{abstract}
Este relatorio apresenta tres propostas de agrupamento aplicadas aos conjuntos
Iris e Wine. A primeira combina o metodo BSAS com K-Means: o BSAS e utilizado
para estimar o numero de agrupamentos por meio da analise da curva
$k(\tau)$, enquanto o K-Means realiza o particionamento final. A segunda
proposta usa Janelas de Parzen para estimar a densidade dos dados e formar
grupos a partir de modos de densidade. A terceira proposta usa KNN por meio de
um grafo de vizinhos compartilhados. Os resultados sao avaliados de forma
qualitativa por visualizacoes em PCA 2D e comparacao com as classes esperadas,
com metricas externas e internas usadas apenas como apoio.
\end{abstract}

\section{Introducao}

O objetivo da atividade e implementar e comparar procedimentos de agrupamento
nao supervisionado em dois conjuntos classicos: Iris Dataset e Wine Dataset.
Embora ambos possuam rotulos conhecidos, esses rotulos nao sao usados no
treinamento dos metodos. Eles servem apenas para avaliacao qualitativa posterior,
permitindo comparar os agrupamentos encontrados com as classes esperadas.

Foram implementadas tres abordagens. Na primeira, o BSAS
(\textit{Basic Sequential Algorithmic Scheme}) estima o numero de grupos
presentes nos dados. Em seguida, esse numero e passado ao K-Means, que executa o
particionamento efetivo. Na segunda abordagem, Janelas de Parzen sao usadas para
estimar uma funcao de densidade, e os grupos sao definidos a partir de regioes
que convergem para o mesmo modo dessa densidade. Na terceira abordagem, um grafo
KNN e construido usando vizinhos compartilhados, e os grupos sao obtidos como
componentes conexas desse grafo.

\section{Pre-processamento}

Antes da aplicacao dos metodos, os dados foram padronizados com
\texttt{StandardScaler}. Assim, cada atributo passa a ter media aproximadamente
zero e desvio-padrao aproximadamente igual a um. Esse pre-processamento e
necessario porque todos os metodos implementados dependem de distancias:
BSAS, K-Means, Janelas de Parzen e KNN.

Sem padronizacao, atributos com maior escala numerica poderiam dominar o calculo
das distancias. Isso seria especialmente problematico no Wine Dataset, pois seus
atributos possuem unidades e amplitudes muito diferentes. Dessa forma, a
padronizacao torna a contribuicao dos atributos mais equilibrada e evita que o
resultado do agrupamento seja determinado artificialmente por uma diferenca de
escala.

Para visualizacao, tambem foi utilizada PCA com duas componentes principais. A
PCA nao foi usada como entrada principal dos algoritmos de agrupamento; ela foi
empregada apenas para projetar os dados em duas dimensoes e permitir uma analise
visual dos grupos encontrados.

\section{Metodologia}

\subsection{BSAS para estimar o numero de grupos}

O BSAS e um metodo sequencial. A cada novo ponto, calcula-se a distancia entre
esse ponto e os representantes dos grupos ja existentes. Se a menor distancia
for menor ou igual a um limiar $\tau$, o ponto e inserido no grupo mais proximo.
Caso contrario, um novo grupo e criado. Assim, $\tau$ controla diretamente o
numero de grupos obtidos.

Valores pequenos de $\tau$ tendem a gerar muitos grupos, pois poucos pontos sao
considerados suficientemente proximos dos grupos existentes. Valores grandes de
$\tau$ tendem a gerar poucos grupos, pois muitos pontos passam a ser aceitos em
grupos ja existentes. No notebook, esse comportamento e estudado pela curva
$k(\tau)$, em que $k$ representa o numero de grupos retornado pelo BSAS para
cada valor de $\tau$.

Como o BSAS e sensivel a ordem de apresentacao dos dados, cada valor de $\tau$
foi avaliado varias vezes com embaralhamentos diferentes. Em seguida, foi usada
a mediana do numero de grupos obtidos. Essa estrategia reduz o efeito de uma
ordem especifica dos dados e torna a estimativa de $k$ mais estavel.

A escolha de $k$ foi feita procurando regioes de estabilidade na curva
$k(\tau)$. Numericamente, foi calculada uma aproximacao da derivada:

\begin{equation}
    \frac{dk}{d\tau} \approx \nabla k(\tau).
\end{equation}

Regioes em que $|dk/d\tau|$ e pequeno indicam que mudancas em $\tau$ alteram
pouco o numero de grupos. Essas regioes foram interpretadas como plateaus de
estabilidade. Por outro lado, regioes com derivada alta indicam transicoes
bruscas, nas quais pequenas variacoes no limiar mudam bastante o agrupamento.

\subsection{K-Means como particionador final}

Apos estimar o numero de grupos com o BSAS, o K-Means foi executado com
\texttt{n\_clusters = k}. Essa combinacao separa dois papeis: o BSAS e usado
como estimador do hiperparametro $k$, e o K-Means e responsavel pelo
particionamento final.

Essa proposta e interessante porque o K-Means normalmente exige que o usuario
defina previamente o numero de grupos. O BSAS fornece uma forma automatizada e
interpretavel de sugerir esse valor a partir da propria estrutura dos dados.

\subsection{Agrupamento por Janelas de Parzen}

Na segunda proposta, foi estimada uma densidade por Janelas de Parzen com kernel
Gaussiano. Para um ponto $x$, a densidade estimada pode ser escrita como:

\begin{equation}
    \hat{p}(x) =
    \frac{1}{n}
    \sum_{i=1}^{n}
    \exp\left(
        -\frac{\lVert x - x_i \rVert^2}{2h^2}
    \right),
\end{equation}

em que $h$ e a largura da janela, tambem chamada de \textit{bandwidth}. A ideia
implementada foi deslocar cada ponto em direcao a regioes de maior densidade,
por meio de uma media ponderada pelos pesos do kernel Gaussiano. Pontos que
convergem para modos proximos da densidade recebem o mesmo rotulo.

Esse procedimento e semelhante em espirito ao Mean Shift, mas foi implementado
explicitamente a partir da estimativa de densidade de Parzen. Alem dos graficos
dos agrupamentos, o notebook tambem exibe um mapa de contorno da densidade
estimada no plano PCA 2D, permitindo observar visualmente os picos de densidade
e sua relacao com os grupos formados.

\subsection{Agrupamento por KNN com vizinhos compartilhados}

Na terceira proposta, foi construido um grafo baseado em KNN. Cada amostra e
representada por um vertice. Dois vertices sao conectados quando satisfazem dois
criterios: eles aparecem como vizinhos mutuos e possuem um numero minimo de
vizinhos compartilhados. Depois disso, os clusters sao definidos como
componentes conexas do grafo.

Essa abordagem e inspirada na ideia de que pontos pertencentes ao mesmo grupo
nao estao apenas proximos, mas tambem compartilham uma vizinhanca local
semelhante. Isso evita algumas conexoes fracas que poderiam surgir quando se usa
apenas a distancia direta entre pares de pontos.

\section{Avaliacao qualitativa}

A avaliacao foi feita principalmente por visualizacao. Para cada dataset, o
notebook gera graficos em PCA 2D comparando:

\begin{itemize}[leftmargin=1.5cm]
    \item as classes esperadas;
    \item os grupos obtidos por BSAS seguido de K-Means;
    \item os grupos obtidos por Janelas de Parzen;
    \item os grupos obtidos pelo metodo KNN com vizinhos compartilhados.
\end{itemize}

Tambem sao exibidas metricas de apoio: ARI, NMI e Silhouette. ARI e NMI comparam
os agrupamentos com os rotulos conhecidos, enquanto a Silhouette avalia a
separacao e coesao dos grupos sem usar os rotulos verdadeiros. Como a proposta
da atividade enfatiza avaliacao qualitativa, essas metricas devem ser lidas
apenas como complemento as visualizacoes.

\section{Discussao dos resultados}

\subsection{Iris Dataset}

No Iris Dataset, espera-se que uma das classes seja visualmente mais separavel,
enquanto as outras duas apresentem sobreposicao parcial. Esse comportamento e
conhecido no conjunto Iris: a classe Setosa costuma ser bem separada, enquanto
Versicolor e Virginica tendem a ocupar regioes mais proximas no espaco de
atributos.

Na abordagem BSAS seguido de K-Means, quando o valor estimado de $k$ se aproxima
do numero esperado de classes, o resultado tende a refletir bem essa estrutura:
um grupo mais isolado e dois grupos com fronteira menos evidente. Entretanto, a
qualidade do resultado depende da estabilidade do plateau escolhido em
$k(\tau)$. Se a curva apresentar varios plateaus plausiveis, a escolha de $k$
pode mudar dependendo da faixa de $\tau$ analisada.

No metodo por Janelas de Parzen, a interpretacao visual da densidade e
particularmente util. Picos de densidade bem separados tendem a corresponder a
grupos mais claros. No caso do Iris, a sobreposicao entre duas classes pode
fazer com que a densidade apresente modos proximos ou ate mesmo um modo
compartilhado, dependendo do \textit{bandwidth}. Assim, o metodo pode fundir
grupos quando a janela e larga demais ou fragmentar grupos quando a janela e
muito estreita.

No metodo KNN com vizinhos compartilhados, a classe mais isolada tende a formar
uma componente bem definida. Ja as classes parcialmente sobrepostas podem ser
ligadas por arestas locais se compartilharem vizinhos suficientes. Portanto, a
fronteira entre os dois grupos menos separados e bastante sensivel aos
parametros $K$ e ao numero minimo de vizinhos compartilhados.

\subsection{Wine Dataset}

No Wine Dataset, a padronizacao e ainda mais importante, pois os atributos
possuem escalas numericas muito diferentes. Apos esse pre-processamento, a PCA
geralmente evidencia uma separacao razoavel entre as classes, embora algumas
regioes possam permanecer proximas.

A combinacao BSAS e K-Means tende a funcionar bem quando a curva $k(\tau)$
apresenta um plateau associado a um numero de grupos compativel com a estrutura
esperada. A vantagem dessa abordagem e a interpretabilidade da escolha de $k$:
o grafico permite visualizar se o numero de grupos foi escolhido em uma regiao
estavel ou em uma transicao abrupta.

No caso do Parzen, a escolha da largura da janela tem impacto direto na
resolucao da densidade. Como Wine possui mais atributos que Iris, a estimativa
de densidade pode ser mais dificil devido ao aumento da dimensionalidade. Em
dimensoes mais altas, os pontos tendem a ficar mais esparsos, o que torna a
escolha do \textit{bandwidth} mais delicada.

O metodo KNN com vizinhos compartilhados pode capturar estruturas locais
interessantes, principalmente quando os grupos possuem densidades diferentes.
Por outro lado, se $K$ for pequeno demais, o grafo pode se fragmentar em muitas
componentes. Se $K$ for grande demais, regioes distintas podem ser conectadas,
reduzindo a separacao entre grupos.

\section{Sensibilidade a parametros}

Os tres procedimentos possuem parametros importantes. A Tabela
\ref{tab:parametros} resume os principais efeitos observados.

\begin{table}[h]
\centering
\caption{Sensibilidade dos metodos aos principais parametros.}
\label{tab:parametros}
\begin{tabular}{lll}
\toprule
Metodo & Parametro & Efeito principal \\
\midrule
BSAS & $\tau$ & Controla a criacao de novos grupos. \\
K-Means & $k$ & Define diretamente o numero de clusters finais. \\
Parzen & $h$ & Controla a suavizacao da densidade estimada. \\
KNN/SNN & $K$ & Controla o tamanho da vizinhanca local. \\
KNN/SNN & vizinhos compartilhados & Controla a rigidez das conexoes do grafo. \\
\bottomrule
\end{tabular}
\end{table}

No BSAS, $\tau$ e o parametro mais critico. Valores baixos levam a muitos
clusters; valores altos levam a poucos clusters. A estrategia de observar
plateaus na curva $k(\tau)$ reduz essa arbitrariedade, mas nao elimina
completamente a necessidade de julgamento.

No K-Means, o parametro $k$ determina o numero final de grupos. Por isso, usar o
BSAS como etapa anterior e uma forma de tornar a escolha mais fundamentada.
Ainda assim, se o BSAS estimar um $k$ inadequado, o K-Means herdara essa
limitacao.

No Parzen, a largura da janela $h$ controla a suavidade da densidade. Um $h$
muito pequeno gera muitos picos locais e pode fragmentar os dados. Um $h$ muito
grande suaviza demais a densidade e pode unir grupos distintos.

No KNN/SNN, $K$ e o numero minimo de vizinhos compartilhados controlam a
conectividade do grafo. Parametros muito permissivos criam conexoes entre grupos
distintos; parametros muito restritivos quebram grupos naturais em varias
componentes pequenas.

\section{Vantagens e desvantagens}

\begin{table}[h]
\centering
\caption{Comparacao qualitativa das abordagens implementadas.}
\label{tab:comparacao}
\begin{tabular}{p{3cm}p{5.4cm}p{5.4cm}}
\toprule
Metodo & Vantagens & Desvantagens \\
\midrule
BSAS + K-Means
& Estima automaticamente $k$; possui interpretacao visual pela curva
$k(\tau)$; K-Means e simples e eficiente.
& BSAS e sensivel a ordem dos dados; K-Means favorece grupos aproximadamente
esfericos; resultado depende do plateau escolhido. \\
\addlinespace
Parzen
& Relaciona agrupamentos a regioes de alta densidade; pode encontrar grupos
sem impor diretamente um valor de $k$; densidade estimada e visualizavel.
& Muito sensivel ao \textit{bandwidth}; sofre em dimensionalidade alta; pode
fundir ou fragmentar grupos dependendo da suavizacao. \\
\addlinespace
KNN/SNN
& Explora estrutura local; pode lidar melhor com formatos nao esfericos; evita
algumas conexoes fracas usando vizinhos compartilhados.
& Sensivel a $K$ e ao criterio de vizinhos compartilhados; pode gerar muitas
componentes pequenas; custo de construir vizinhancas cresce com o numero de
amostras. \\
\bottomrule
\end{tabular}
\end{table}

\section{Limitacoes}

Uma limitacao geral e que as visualizacoes em PCA 2D mostram apenas uma
projecao dos dados. Separacoes ou sobreposicoes vistas no grafico podem nao
representar perfeitamente a geometria no espaco original. Portanto, os graficos
devem ser interpretados como apoio qualitativo, nao como prova absoluta da
estrutura dos grupos.

Outra limitacao e a dependencia de parametros. Embora o notebook proponha regras
automaticas ou semi-automaticas para definir alguns valores, ainda existe
sensibilidade a escolhas como $\tau$, \textit{bandwidth}, $K$ e numero minimo de
vizinhos compartilhados. Uma analise mais completa poderia testar grades de
parametros e avaliar a estabilidade dos agrupamentos.

O BSAS tambem apresenta dependencia da ordem dos dados. A implementacao reduz
esse problema ao executar multiplos embaralhamentos e usar a mediana de $k$, mas
isso nao remove totalmente a caracteristica sequencial do metodo.

Por fim, os datasets Iris e Wine sao relativamente pequenos e bem conhecidos.
Em bases maiores, com ruido, atributos irrelevantes ou estruturas mais
complexas, o comportamento dos metodos pode ser diferente. Especialmente no
caso de Parzen, a estimativa de densidade pode se tornar mais dificil em
dimensoes elevadas.

\section{Conclusao}

A atividade mostrou que diferentes ideias de agrupamento enfatizam aspectos
distintos dos dados. A combinacao BSAS + K-Means e uma alternativa util quando
se deseja estimar o numero de clusters antes de aplicar um particionador
classico. A analise da curva $k(\tau)$ e de sua derivada numerica fornece uma
justificativa visual e quantitativa para a escolha de $k$.

O metodo baseado em Janelas de Parzen oferece uma interpretacao por densidade:
grupos sao associados a regioes de maior concentracao de pontos. Essa abordagem
e intuitiva e visualmente rica, mas depende fortemente da largura da janela. O
metodo KNN/SNN, por sua vez, modela a estrutura local dos dados por meio de um
grafo, sendo interessante para identificar grupos conectados por vizinhancas
semelhantes.

Em termos praticos, nenhuma abordagem e universalmente superior. BSAS + K-Means
e simples e eficiente, mas assume uma estrutura de grupos mais regular. Parzen e
mais flexivel em termos de densidade, mas e sensivel a suavizacao. KNN/SNN
captura relacoes locais, mas depende da conectividade do grafo. Assim, a melhor
interpretacao dos resultados surge da combinacao entre visualizacao, metricas de
apoio e analise critica da sensibilidade dos parametros.

\end{document}
